\ifdefined\gkver\else\def\gkver{student}\fi
\documentclass[\gkver]{gaokaozhenti}
\begin{document}

\chapter{2007年北京理综}
%% paperType: 理综物理部分

\begin{enumerate}
\item
%% number: 13
%% paperName: 2007年北京理综第 13 题 6 分
%% typeId: 1
%% type: 单选题
%% chapter: 光学
%% point: 全反射
%% method: 无
%% score: 6
%% degree: 650
%% duplicateId: 0
%% body:
光导纤维的结构如图所示，其内芯和外套材料不同，光在内芯中传播。以下关于光导纤维的说法正确的是\xzanswer{A}

\onepicture[w=6.0cm]{\includesvg{figs/13}}

\fourchoices[answer=A]
{内芯的折射率比外套大，光传播时在内芯与外套的界面发生全反射}
{内芯的折射率比外套小，光传播时在内芯与外套的界面发生全反射}
{内芯的折射率比外套小，光传播时在内芯与外套的界面发生折射}
{内芯的折射率比外套相同，外套的材料有韧性，可以起保护作用}

%% answer: A
%% memo:
\memoanswer{本题原卷及材料中未提供详解，待补充。}

\item
%% number: 14
%% paperName: 2007年北京理综第 14 题 6 分
%% typeId: 1
%% type: 单选题
%% chapter: 原子和原子核
%% point: 玻尔模型
%% method: 无
%% score: 6
%% degree: 650
%% duplicateId: 0
%% body:
下列说法正确的是\xzanswer{D}

\fourchoices[answer=D]
{太阳辐射的能量主要来自太阳内部的核裂变反应}
{汤姆生发现电子，表明原子具有核式结构}
{一束光照射到某种金属上不能发生光电效应，是因为该束光的波长太短}
{按照波尔理论，氢原子核外电子从半径较小的轨道跃迁到半径较大的轨道时，电子的动能减小，原子总能量增加}

%% answer: D
%% memo:
\memoanswer{本题原卷及材料中未提供详解，待补充。}

\item
%% number: 15
%% paperName: 2007年北京理综第 15 题 6 分
%% typeId: 1
%% type: 单选题
%% chapter: 曲线运动
%% point: 人造地球卫星
%% method: 无
%% score: 6
%% degree: 550
%% duplicateId: 0
%% body:
不久前欧洲天文学就发现了一颗可能适合人类居住的行星，命名为“格利斯581c”。该行星的质量是地球的5倍，直径是地球的1.5倍。设想在该行星表面附近绕行星沿圆轨道运行的人造卫星的动能为 $E_{k1}$，在地球表面附近绕地球沿圆轨道运行的相同质量的人造卫星的动能为 $E_{k2}$，则 $E_{k1}/E_{k2}$ 为\xzanswer{C}

\fourchoices[answer=C]
{0.13}
{0.3}
{3.33}
{7.5}

%% answer: C
%% memo:
\memoanswer{本题原卷及材料中未提供详解，待补充。}

\item
%% number: 16
%% paperName: 2007年北京理综第 16 题 6 分
%% typeId: 1
%% type: 单选题
%% chapter: 其他
%% point: 物理思想
%% method: 控制变量法
%% score: 6
%% degree: 650
%% duplicateId: 0
%% body:
为研究影响家用保温瓶保温效果的因素，某同学在保温瓶中灌入热水，现测量初始水温，经过一段时间后再测量末态水温。改变实验条件，先后共做了6次实验，实验数据记录如下表：

\begin{center}
\begin{tabular}{|c|c|c|c|c|}
\hline
序号 & 瓶内水量（$\UmL$） & 初始水温（$\UdC$） & 时间（$\Uh$） & 末态水温（$\UdC$）\\
\hline
1 & 1000 & 91 & 4 & 78\\
\hline
2 & 1000 & 98 & 8 & 74\\
\hline
3 & 1500 & 91 & 4 & 80\\
\hline
4 & 1500 & 98 & 10 & 75\\
\hline
5 & 2000 & 91 & 4 & 82\\
\hline
6 & 2000 & 98 & 12 & 77\\
\hline
\end{tabular}
\end{center}

下列研究方案中符合控制变量方法的是\xzanswer{A}

\fourchoices[answer=A]
{若研究瓶内水量与保温效果的关系，可用第1、3、5次实验数据}
{若研究瓶内水量与保温效果的关系，可用第2、4、6次实验数据}
{若研究初始水温与保温效果的关系，可用第1、2、3次实验数据}
{若研究保温时间与保温效果的关系，可用第4、5、6次实验数据}

%% answer: A
%% memo:
\memoanswer{本题原卷及材料中未提供详解，待补充。}

\item
%% number: 17
%% paperName: 2007年北京理综第 17 题 6 分
%% typeId: 1
%% type: 单选题
%% chapter: 交流电
%% point: 交流电
%% method: 无
%% score: 6
%% degree: 650
%% duplicateId: 0
%% body:
电阻 $R_{1}$、$R_{2}$ 交流电源按照图 1 所示方式连接，$R_{1}=10\UO$，$R_{2}=20\UO$。合上开关 S 后，通过电阻 $R_{2}$ 的正弦交变电流 $i$ 随时间 $t$ 变化的情况如图 2 所示。则\xzanswer{B}

\twopicture[labelprefix={图}, numA=1, numB=2, labelA=2007北京理综17a, labelB=2007北京理综17b, wA=4.0cm, wB=6.2cm]
{figs/17a.png}
{figs/17b.png}

\fourchoices[answer=B]
{通过 $R_{1}$ 的电流的有效值是 $1.2\UA$}
{$R_{1}$ 两端的电压有效值是 $6\UV$}
{通过 $R_{2}$ 的电流的有效值是 $1.2\sqrt{2}\UA$}
{$R_{2}$ 两端的电压有效值是 $6\sqrt{2}\UV$}

%% answer: B
%% memo:
\memoanswer{本题原卷及材料中未提供详解，待补充。}

\item
%% number: 18
%% paperName: 2007年北京理综第 18 题 6 分
%% typeId: 1
%% type: 单选题
%% chapter: 直线运动
%% point: 匀速直线运动
%% method: 近似与估算
%% score: 6
%% degree: 650
%% duplicateId: 0
%% body:
图示为高速摄影机拍摄到的子弹穿过苹果瞬间的照片。该照片经过放大后分析出，在曝光时间内，子弹影像前后错开的距离约为子弹长度的 $1\%\sim2\%$。已知子弹飞行速度约为 $500\Ums$，因此可估算出这幅照片的曝光时间最接近\xzanswer{B}

\onepicture[w=6.0cm]{figs/18.jpg}

\fourchoices[answer=B]
{$10^{-3}\Us$}
{$10^{-6}\Us$}
{$10^{-9}\Us$}
{$10^{-12}\Us$}

%% answer: B
%% memo:
\memoanswer{本题原卷及材料中未提供详解，待补充。}

\item
%% number: 19
%% paperName: 2007年北京理综第 19 题 6 分
%% typeId: 1
%% type: 单选题
%% chapter: 机械振动
%% point: 单摆
%% method: 无
%% score: 6
%% degree: 650
%% duplicateId: 0
%% body:
如图所示的单摆，摆球 a 向右摆动到最低点时，恰好与一沿水平方向向左运动的粘性小球 b 发生碰撞，并粘在一起，且摆动平面不变。已知碰撞前 a 球摆动的最高点与最低点的高度差为 $h$，摆动的周期为 $T$，a 球质量是 b 球质量的 5 倍，碰撞前 a 球在最低点的速度是 b 球速度的一半。则碰撞后\xzanswer{D}

\onepicture[w=3.5cm]{figs/19.png}

\fourchoices[answer=D]
{摆动的周期为 $\sqrt{\frac{5}{6}}T$}
{摆动的周期为 $\sqrt{\frac{6}{5}}T$}
{摆球最高点与最低点的高度差为 $0.3h$}
{摆球最高点与最低点的高度差为 $0.25h$}

%% answer: D
%% memo:
\memoanswer{本题原卷及材料中未提供详解，待补充。}

\item
%% number: 20
%% paperName: 2007年北京理综第 20 题 6 分
%% typeId: 1
%% type: 单选题
%% chapter: 功和能
%% point: 动能定理
%% method: 无
%% score: 6
%% degree: 450
%% duplicateId: 0
%% body:
在真空中的光滑水平绝缘面上有一带电小滑块。开始时滑块静止。若在滑块所在空间加一水平匀强电场 $E_{1}$，持续一段时间后立即换成与 $E_{1}$ 相反方向的匀强电场 $E_{2}$。当电场 $E_{2}$ 与电场 $E_{1}$ 持续时间相同时，滑块恰好回到初始位置，且具有动能 $E_{k}$。在上述过程中，$E_{1}$ 对滑块的电场力做功为 $W_{1}$，冲量大小为 $I_{1}$；$E_{2}$ 对滑块的电场力做功为 $W_{2}$，冲量大小为 $I_{2}$。则\xzanswer{C}

\fourchoices[answer=C]
{$I_{1}=I_{2}$}
{$4I_{1}=I_{2}$}
{$W_{1}=0.25E_{k}$，$W_{2}=0.75E_{k}$}
{$W_{1}=0.20E_{k}$，$W_{2}=0.80E_{k}$}

%% answer: C
%% memo:
\memoanswer{本题原卷及材料中未提供详解，待补充。}

\item
%% number: 21
%% paperName: 2007年北京理综第 21 题 15 分
%% typeId: 4
%% type: 实验题
%% chapter: 运动和力
%% point: 牛顿第二定律
%% method: 无
%% score: 15
%% degree: 450
%% duplicateId: 0
%% body:
\begin{enumerate}
	\item
	如图所示是电子射线管的示意图。接通电源后，电子射线由阴极沿 $x$ 轴方向射出，在荧光屏上会看到一条亮线。要使荧光屏上的亮线向下（$z$ 轴方向）偏转，在下列措施中可采用的是 \tkanswer[1.5cm]{B}（填选项代号）。

	\onepicture[w=6.0cm]{figs/21.png}

	\fourchoices[answer=B]
	{加一磁场，磁场方向沿 $z$ 轴负方向}
	{加一磁场，磁场方向沿 $y$ 轴正方向}
	{加一电场，电场方向沿 $z$ 轴负方向}
	{加一电场，电场方向沿 $y$ 轴正方向}

	\item
	某同学用如图所示的实验装置研究小车在斜面上的运动。实验步骤如下：

	\begin{enumerate}[label=(\Alph*)]
		\item
		安装好实验器材。
		\item
		接通电源后，让拖着纸带的小车沿平板斜面向下运动，重复几次。选出一条点迹比较清晰的纸带，舍去开始密集的点迹，从便于测量的点开始，每两个打点间隔取一个计数点，如图中 0、1、2……6 点所示。
		\item
		测量 1、2、3……6 计数点到 0 计数点的距离，分别记作：$s_{1}$、$s_{2}$、$s_{3}$……$s_{6}$。
		\item
		通过测量和计算，该同学判断出小车沿平板做匀加速直线运动。
		\item
		分别计算出 $s_{1}$、$s_{2}$、$s_{3}$……$s_{6}$ 与对应时间的比值 $\frac{s_{1}}{t_{1}}$、$\frac{s_{2}}{t_{2}}$、$\frac{s_{3}}{t_{3}}$……$\frac{s_{6}}{t_{6}}$。
		\item
		以 $\frac{s}{t}$ 为纵坐标、$t$ 为横坐标，标出 $\frac{s}{t}$ 与对应时间 $t$ 的坐标点，画出 $\frac{s}{t}-t$ 图线。
	\end{enumerate}

	结合上述实验步骤，请你完成下列任务：

	\begin{steps}
		\item
		实验中，除打点计时器（含纸带、复写纸）、小车、平板、铁架台、导线及开关外，在下面的仪器和器材中，必须使用的有 \tkanswer[1cm]{A} 和 \tkanswer[1cm]{C}。（填选项代号）

		（A）电压合适的 $50\UHz$ 交流电源　　（B）电压可调的直流电源

		（C）刻度尺　　（D）秒表　　（E）天平　　（F）重锤
		\item
		将最小刻度为 $1\Umm$ 的刻度尺的 0 刻线与 0 计数点对齐，0、1、2、5 计数点所在位置如图所示，则 $s_{2}=$ \tkanswer[2cm]{2.97$\sim$2.99}$\Ucm$，$s_{5}=$ \tkanswer[2cm]{13.19$\sim$13.21}$\Ucm$。
		\item
		该同学在图中已标出 1、3、4、6 计数点对应的坐标，请你在该图中标出与 2、5 两个计数点对应的坐标点，并画出 $\frac{s}{t}-t$。
		\item
		根据 $\frac{s}{t}-t$ 图线判断，在打 0 计数点时，小车的速度 $v_{0}=$ \tkanswer[2cm]{0.16$\sim$0.20}$\Ums$；它在斜面上运动的加速度 $a=$ \tkanswer[2cm]{4.50$\sim$5.10}$\Umsq$。
	\end{steps}
\end{enumerate}

\fourpicture[labelA=2007北京理综21a, labelB=2007北京理综21b, labelC=2007北京理综21c, labelD=2007北京理综21d, align=right, wA=4.5cm, wB=8cm, wC=8.5cm, wD=4cm]
{figs/21a.png}
{figs/21b.png}
{figs/21c.png}
{figs/21d.png}

%% answer: 
\jdanswer{
\begin{enumerate}
	\item
	B
	\item
	① A、C；② 2.97$\sim$2.99，13.19$\sim$13.21；③ 略；④ 0.16$\sim$0.20，4.50$\sim$5.10。
\end{enumerate}
}
%% memo:
\memoanswer{本题原卷及材料中未提供详解，待补充。}

\item
%% number: 22
%% paperName: 2007年北京理综第 22 题 16 分
%% typeId: 6
%% type: 计算题
%% chapter: 电场
%% point: 带电粒子的偏转
%% method: 无
%% score: 16
%% degree: 450
%% duplicateId: 0
%% body:
两个半径均为 $R$ 的圆形平板电极，平行正对放置，相距为 $d$，极板间的电势差为 $U$，板间电场可以认为是均匀的。一个 $\alpha$ 粒子从正极板边缘以某一初速度垂直于电场方向射入两极板之间，到达负极板时恰好落在极板中心。已知质子电荷为 $e$，质子和中子的质量均视为 $m$，忽略重力和空气阻力的影响，求：

\begin{enumerate}
	\item
	极板间的电场强度 $E$；
	\item
	$\alpha$ 粒子在极板间运动的加速度 $a$；
	\item
	$\alpha$ 粒子的初速度 $v_{0}$。
\end{enumerate}

%% answer: 
\jdanswer{
\begin{enumerate}
	\item
	$E=\frac{U}{d}$
	\item
	$a=\frac{eU}{2md}$
	\item
	$v_{0}=\frac{R}{2d}\sqrt{\frac{eU}{m}}$
\end{enumerate}
}
%% memo:
\memoanswer{
（1）极板间的电场强度 $E=\frac{U}{d}$。

（2）$\alpha$ 粒子带两个单位正电荷，电荷量 $q=2e$，质量 $m'=4m$。由牛顿第二定律，$\alpha$ 粒子在极板间运动的加速度 $a=\frac{qE}{m'}=\frac{2eU}{4md}=\frac{eU}{2md}$。

（3）$\alpha$ 粒子垂直于电场方向射入，在极板间做类平抛运动。沿电场方向有 $d=\frac{1}{2}at^{2}$，得 $t=\sqrt{\frac{2d}{a}}=2d\sqrt{\frac{m}{eU}}$；沿初速度方向有 $R=v_{0}t$，所以 $v_{0}=\frac{R}{t}=\frac{R}{2d}\sqrt{\frac{eU}{m}}$。
}

\item
%% number: 23
%% paperName: 2007年北京理综第 23 题 18 分
%% typeId: 6
%% type: 计算题
%% chapter: 功和能
%% point: 交通工具启动
%% method: 无
%% score: 18
%% degree: 450
%% duplicateId: 0
%% body:
环保汽车将为 2008 年奥运会场馆服务。某辆以蓄电池为驱动能源的环保汽车，总质量 $m=3\times10^{3}\Ukg$。当它在水平路面上以 $v=36\Ukmh$ 的速度匀速行驶时，驱动电机的输入电流 $I=50\UA$，电压 $U=300\UV$。在此行驶状态下

\begin{enumerate}
	\item
	求驱动电机的输入功率 $P_{\text{电}}$；
	\item
	若驱动电机能够将输入功率的 $90\%$ 转化为用于牵引汽车前进的机械功率 $P_{\text{机}}$，求汽车所受阻力与车重的比值（$g$ 取 $10\Umsq$）；
	\item
	设想改用太阳能电池给该车供电，其他条件不变，求所需的太阳能电池板的最小面积。结合计算结果，简述你对该设想的思考。已知太阳辐射的总功率 $P_{0}=4\times10^{26}\UW$，太阳到地球的距离 $r=1.5\times10^{11}\Um$，太阳光传播到达地面的过程中大约有 $30\%$ 的能量损耗，该车所用太阳能电池的能量转化效率约为 $15\%$。
\end{enumerate}

%% answer: 
\jdanswer{
\begin{enumerate}
	\item
	$P_{\text{电}}=1.5\times10^{4}\UW$
	\item
	$0.045$
	\item
	$S=101\Umq$
\end{enumerate}
}
%% memo:
\memoanswer{
（1）驱动电机的输入功率 $P_{\text{电}}=IU=50\times300\UW=1.5\times10^{4}\UW$。

（2）汽车匀速行驶时，$P_{\text{机}}=0.9P_{\text{电}}=fv$，其中 $v=36\Ukmh=10\Ums$，所以阻力 $f=\frac{0.9P_{\text{电}}}{v}=\frac{0.9\times1.5\times10^{4}}{10}\UN=1.35\times10^{3}\UN$，阻力与车重之比 $\frac{f}{mg}=\frac{1.35\times10^{3}}{3\times10^{3}\times10}=0.045$。

（3）当阳光垂直电池板入射时，所需电池板面积最小，设最小面积为 $S$，距太阳中心 $r$ 的球面面积 $S_{0}=4\pi r^{2}$。设电池板实际接收到的太阳能功率为 $P$，若不考虑能量损耗，电池板接收的太阳能功率为 $P'$，则 $\frac{P'}{P_{0}}=\frac{S}{S_{0}}$；实际 $P=(1-30\%)P'$，故 $\frac{P}{P_{0}(1-30\%)}=\frac{S}{S_{0}}$。又 $P_{\text{电}}=15\%P$，所以电池板的最小面积 $S=\frac{PS_{0}}{0.7P_{0}}=\frac{4\pi r^{2}P_{\text{电}}}{0.15\times0.7P_{0}}=101\Umq$。现在还不能达到设计要求，要进一步提高太阳能电池的转化率，减小车的质量，提高电动机的效率。
}

\item
%% number: 24
%% paperName: 2007年北京理综第 24 题 20 分
%% typeId: 6
%% type: 计算题
%% chapter: 电磁感应
%% point: 电磁感应能量分析
%% method: 无
%% score: 20
%% degree: 450
%% duplicateId: 0
%% body:
用密度为 $D$、电阻率为 $\rho$、横截面积为 $A$ 的薄金属条制成边长为 $L$ 的闭合正方形框 abb$'$a$'$。如图所示，金属方框水平放在磁极的狭缝间，方框平面与磁场方向平行。设匀强磁场仅存在于相对磁极之间，其他地方的磁场忽略不计。可认为方框的 aa$'$ 边和 bb$'$ 边都处在磁极之间，极间磁感应强度大小为 $B$。方框从静止开始释放，其平面在下落过程中保持水平（不计空气阻力）。

\begin{enumerate}
	\item
	求方框下落的最大速度 $v_{m}$（设磁场区域在竖直方向足够长）；
	\item
	当方框下落的加速度为 $g/2$ 时，求方框的发热功率 $P$；
	\item
	已知方框下落时间为 $t$ 时，下落高度为 $h$，其速度为 $v_{t}$（$v_{t}<v_{m}$）。若在同一时间 $t$ 内，方框内产生的热与一恒定电流 $I_{0}$ 在该框内产生的热相同，求恒定电流 $I_{0}$ 的表达式。
\end{enumerate}

\twopicture[labelprefix={图}, numA=1, numB=2, labelA=2007北京理综24a, labelB=2007北京理综24b, wA=5.0cm, wB=5.5cm, align=right]
{figs/24a.png}
{figs/24b.png}

%% answer: 
\jdanswer{
\begin{enumerate}
	\item
	$v_{m}=\frac{4\rho Dg}{B^{2}}$
	\item
	$P=\frac{4\rho ALD^{2}g^{2}}{B^{2}}$
	\item
	$I_{0}=A\sqrt{\frac{D}{2\rho t}(2gh-v_{t}^{2})}$
\end{enumerate}
}
%% memo:
\memoanswer{
（1）方框质量 $m=4LAD$，方框电阻 $R=\frac{4\rho L}{A}$。方框下落速度为 $v$ 时，产生的感应电动势 $E=2BLv$，感应电流 $I=\frac{E}{R}=\frac{BAv}{2\rho}$，所受安培力 $F=2BIL=\frac{B^{2}ALv}{\rho}$，方向竖直向上。当 $F=mg$ 时方框达到最大速度 $v_{m}$，即 $\frac{B^{2}ALv_{m}}{\rho}=4LADg$，解得 $v_{m}=\frac{4\rho Dg}{B^{2}}$。

（2）方框下落加速度为 $\frac{g}{2}$ 时，由牛顿第二定律 $mg-2BIL=m\cdot\frac{g}{2}$，得 $I=\frac{mg}{4BL}=\frac{ADg}{B}$，方框的发热功率 $P=I^{2}R=\frac{4\rho ALD^{2}g^{2}}{B^{2}}$。

（3）根据能量守恒定律 $mgh=\frac{1}{2}mv_{t}^{2}+I_{0}^{2}Rt$，得 $I_{0}=\sqrt{\frac{m}{Rt}\left(gh-\frac{1}{2}v_{t}^{2}\right)}$；代入 $m=4LAD$、$R=\frac{4\rho L}{A}$，解得 $I_{0}=A\sqrt{\frac{D}{2\rho t}(2gh-v_{t}^{2})}$。
}

\end{enumerate}

\end{document}
